% TOP_PROBLEM: {"id": 3338, "title": "Order-invariant queries", "category_id": 18, "category": {"id": 18, "name": "logic", "display_name": "Logic", "description": "Problems in mathematical logic, model theory, proof theory, and finite model theory.", "slug": "logic", "order_index": 18, "created_at": "2026-07-31T15:26:25.671Z"}, "status": "open"}
% FIRST_POSTED: null
% ENTRY_KIND: statement_only
% Imported from: batch03.tex; UnsolvedMath ID 3338
\documentclass[11pt]{article}
\usepackage[margin=25mm]{geometry}
\usepackage{fontspec}
\setmainfont{Latin Modern Roman}
\usepackage{amsmath,amssymb,mathtools}
\usepackage{unicode-math}
\setmathfont{Latin Modern Math}
\usepackage{xurl,hyperref}
\hypersetup{colorlinks=true,urlcolor=blue}
\setcounter{secnumdepth}{0}
\setlength{\parindent}{0pt}
\setlength{\parskip}{5pt}
\setlength{\emergencystretch}{3em}
\newcommand{\sref}[2]{\hyperlink{source:#1:#2}{[S#2]}}
\newcommand{\cataloguescope}[1]{\paragraph{Recorded target.}#1}
\begin{document}
% UnsolvedMath ID 3338; source catalogue code OPG-37429
\setcounter{section}{3337}
\section{Order-invariant queries}\label{3338}
{\small\sffamily UnsolvedMath ID 3338 · Logic}
\subsection{Definitions and mathematical statement}

\paragraph{Shared notation.}$\mathbb N=\{1,2,\ldots\}$, and $\mathbb Z,\mathbb Q,\mathbb R,\mathbb C$ denote the integers, rationals, reals and complex numbers. Any other conventions are specified locally.


Work with finite graphs carrying a fixed finite vocabulary of unary vertex-colour predicates. First-order logic (FO) quantifies over vertices; monadic second-order logic (MSO) also quantifies over vertex subsets. A formula $\varphi$ using an additional linear order $<$ is order-invariant if its truth, including for each assignment to free variables, is independent of the chosen vertex order. Write $\operatorname{InvFO}(<)$ for the resulting queries on unordered graphs. Expressive-power equality or inclusion means equivalence to a formula of the other logic on the specified graph class. A tree decomposition consists of vertex bags on a tree, covering edges and satisfying connectedness of the bags containing each vertex; tree-width is the least maximum bag size minus one.\par
A query is Gaifman-local if a fixed-radius induced neighborhood of its parameter tuple determines its value within a fixed structure. It is Hanf-local if, for some fixed radius, agreement of the counts of rooted neighborhood isomorphism types (with the parameters distinguished) forces agreement of the query value. A $0$--$1$ law means that each sentence has limiting probability zero or one in the usual uniform random labeled finite-structure model. Effective syntax means a computably recognizable collection of formulas with effective finite-structure semantics.\par
For addition-invariant FO, label the vertices bijectively by an initial segment of the positive integers and add the induced ternary relation $x+y=z$; require answers to be independent of that bijection. On words, the underlying word order and letter predicates remain part of the input, while the extra numerical labeling is arbitrary. A regular word language is one accepted by a finite automaton.
\paragraph{Statement recorded in the supplied source.}
The source asks all of the following. (1) For every fixed tree-width bound $b$, does $\operatorname{InvFO}(<)=\mathrm{FO}$ hold on graphs of tree-width at most $b$? (2) Is $\operatorname{InvFO}(<)\subseteq\mathrm{MSO}$ on all finite graphs? (3) Does order-invariant FO have a $0$--$1$ law? (4) Is every order-invariant FO query Hanf-local? (5) Is there a logic with effective syntax capturing exactly order-invariant FO? Its additional questions are (6) is addition-invariant FO Gaifman-local, and (7) can addition-invariant FO define a nonregular language of finite words? \sref{3338}{2}
\subsection{Short English statement}
The seven questions concern whether order-invariant first-order logic collapses on bounded-tree-width graphs, lies in monadic second-order logic, has a zero-one law or locality, and admits effective syntax, together with locality and word-language expressiveness for addition-invariant logic.
\subsection{Sources}
\begin{description}
\item[\hypertarget{source:3338:1}{S1}] ULAM.AI and the UnsolvedMath Contributors, UnsolvedMath: A Curated Collection of Open Mathematics Problems, record 3338 (OPG-37429). CC BY 4.0. \url{https://huggingface.co/datasets/ulamai/UnsolvedMath}.
\item[\hypertarget{source:3338:2}{S2}] Open Problem Garden, Order-invariant queries, node 37429, original problem page. Original question and full discussion. \url{https://www.openproblemgarden.org/op/order_invariant_queries}.

\end{description}
\label{end:3338}
\end{document}
